% ============================================================
% 00_abstract.tex — Abstract (MK + EN)
% ============================================================

\section*{Апстракт на македонски јазик}

Оваа дипломска работа го претставува дизајнот, имплементацијата и евалуацијата на целосно автоматизиран систем за управување со состав во играта Fantasy Premier League (FPL), изграден врз трослојна хибридна архитектура на вештачка интелигенција. Системот ги обединува машинското учење, комбинаторната оптимизација и големите јазични модели (LLM) со цел самостојно да донесува одлуки за избор на состав, трансфери, избор на капитен и активирање на чипови во текот на целата сезона, без рачна интервенција.

Слојот за предвидување користи четири одделни LightGBM модели — по еден за секоја позиција (голман, одбранбен играч, везист, напаѓач) — за да ги предвиди очекуваните поени за секоја натпреварувачка недела. Моделите се тренирани врз шест сезони историски FPL податоци (2019--20 до 2024--25), вкупно 51\,044 редови набљудувања на ниво на играч и натпреварувачка недела, организирани во четири позициски датотеки. Карактеристиките опфаќаат прозорци на моментна форма (последните 3 и 5 натпреварувачки недели), очекувани голови и очекувани примени голови на ниво на тим (xG/xGA) од Understat, оцени за тежина на натпреварите (FDR), пазарни сигнали (обем на трансфери, процент на сопственост) и статистики за нови трансфери од девет европски лиги преземени од FBref. Тренирањето следи протокол на прошетна-напред крос-валидација со пет временски подредени набори и експоненцијално растечки тежини, со цел да им се даде поголема важност на поновите сезони.

Слојот за оптимизација го формулира изборот на состав како проблем на целобројно линеарно програмирање (ILP), кој ја максимизира сумата на предвидени поени за состав од 15 играчи во рамки на FPL ограничувањата: буџет од \pounds 100 милиони, барања за позициска структура, најмногу тројца играчи од еден клуб и казни за трансфер од $-4$ поени за секој дополнителен трансфер над дозволениот.

Слојот на интелигенција се состои од седум модули кои собираат информации во реално време пред секој рок: статус на повреди и достапност на играчите од FPL API, транскрипти од прес-конференции преземени од Fantasy Football Scout, оцени за достапност по играч (0--100), оцени за ризик од ротација (0--100) и стратешки препораки генерирани со LLM преку Gemini 2.5 Flash API. Овие сигнали се претвораат во мултипликативни корективни фактори кои се применуваат врз предвидувањата на моделот пред оптимизацијата, со што директно се казнуваат изборите со зголемен ризик од повреда или ротација.

Системот беше евалуиран на натпреварувачките недели 1--28 од сезоната 2025--26 во FPL, при што постигна вкупно \textbf{1\,799 поени} (просек од \textbf{64,3} поени по натпреварувачка недела), без ниту еден казнен поен од трансфери. Хиперпараметрите беа подесени преку заедничко случајно пребарување од 250 обиди, проследено со Бајесова оптимизација од 100 обиди преку TPE семплерот на Optuna, при што финалната конфигурација (Optuna обид 429) ги надмина сите претходни конфигурации. Резултатите покажуваат дека интелигентниот цевковод додава значителна вредност над чистото предвидување засновано на машинско учење и дека LightGBM конзистентно го надминува XGBoost низ сите истражени конфигурации на хиперпараметри — со 14 од 15 најдобри резултати кои му припаѓаат на LightGBM.

\bigskip
\noindent\textbf{Клучни зборови:} машинско учење, оптимизација, фантази спорт, LightGBM, ILP, LLM, LLM агент, Бајесова оптимизација, временски редови, Fantasy Premier League

